\section{用户线程怎样等待同一把锁}

Process 是地址空间、文件表和目录上下文的资源容器；Thread 才是处理器可以
选择、阻塞和恢复的执行流。增加用户线程不是复制一个栈指针，而是同时建立
调度身份、两级栈、完整处理器现场、TLS、等待关系、退出值和可回滚所有权。

\topic{历史动机：共享资源为什么不等于共享执行现场}
早期 Unix 主要用进程表达并发。每个进程有独立地址空间，隔离清晰，但两个
协作执行流若只想共享一张缓存表，仍要经 pipe、共享内存或文件传递。图形界面、
数据库和网络服务器逐渐需要在同一地址空间中安排多个阻塞任务：代码、堆、
打开文件和当前目录应该共享，RIP、RSP、寄存器和等待原因却必须独立。

这形成今天仍重要的分界：

\begin{center}
\begin{osgridtabular}{L{0.27\linewidth}L{0.27\linewidth}L{0.34\linewidth}}
事实 & 所有者 & 不共享的原因或共享收益 \\ \hline
页表根、VMA 与后备 & Process & 所有 Thread 观察同一组用户对象 \\ \hline
FileTable 与 FsContext & Process & fd、共享 offset 和 cwd 形成同一环境 \\ \hline
PID 与父子关系 & Process & wait 消费的是程序实例的终止记录 \\ \hline
RIP/RSP/RFLAGS & Thread & 每个执行流停在不同指令和调用深度 \\ \hline
通用/FXSAVE 现场 & Thread & 整数、x87、MMX、XMM 状态不能串用 \\ \hline
内核栈与用户栈 & Thread & 特权入口和普通调用都需要独立生存区间 \\ \hline
TID、等待节点与退出值 & Thread & 调度、join 和同步以执行流为单位 \\ \hline
FS-base TLS & Thread & 同一条机器指令要定位不同私有实例 \\ \hline
\end{osgridtabular}
\end{center}

Linux 后来用可组合的共享标志构造任务；POSIX pthread 又为应用提供可移植
接口。本项目不先复制两套庞大兼容层，而是在已经分离的 Process/Thread
内核模型上定义最小公开 ABI。这样每个新增字段都有硬件或所有权原因，而不是
为了让接口名字看起来像现成系统。

\topic{一条 Thread 从用户请求到 Ready 发布}
用户运行时先保留 \(68\,\mathrm{KiB}\) 栈区：低端一页撤销 VMA，成为永久
not-present guard；其上 \(64\,\mathrm{KiB}\) 为读写、不可执行栈。TLS 另占
一张读写、不可执行页。两处都来自现有匿名 VMA，而不是 Kernel 私自选择用户
地址。

\begin{pdfdiagram}[node distance=7mm and 8mm]
  \node[os block] (map) {用户 MapAnonymous\\栈保留 + TLS};
  \node[os state,right=of map] (guard) {撤销低端 4 KiB\\形成 guard};
  \node[os block,right=of guard] (request) {48 B\\ThreadCreateRequest};
  \node[os hardware,below=of request] (validate) {Kernel 验证\\VMA/权限/对齐};
  \node[os state,left=of validate] (prepare) {KernelStack\\UserContext/FXSAVE};
  \node[os block,left=of prepare] (publish) {最后发布 Ready\\允许 PIT 选择};
  \draw[os arrow] (map) -- (guard);
  \draw[os arrow] (guard) -- (request);
  \draw[os arrow] (request) -- (validate);
  \draw[os arrow] (validate) -- (prepare);
  \draw[os arrow] (prepare) -- (publish);
\end{pdfdiagram}
\diagramnote{“最后发布”是并发正确性边界：Ready 不是初始化中的中间状态。}

公开请求只含固定宽度字段：

\begin{center}
\begin{osgridtabular}{L{0.30\linewidth}L{0.19\linewidth}L{0.39\linewidth}}
字段 & 宽度 & Kernel 要证明的性质 \\ \hline
entry address & 64 bit & 用户规范地址，位于可执行且不可写 VMA \\ \hline
argument & 64 bit & 原样放入初始 RDI，不由 Kernel 解引用 \\ \hline
stack base/size & 各 64 bit & 页对齐、至少 16 KiB、读写且 NX、无溢出 \\ \hline
stack pointer & 64 bit & 精确位于栈顶减 8，满足入口 ABI \\ \hline
TLS base & 64 bit & 16 字节对齐，当前地址空间内可写 \\ \hline
\end{osgridtabular}
\end{center}

初始 RSP 选择为

\[
  RSP_0 = stack_{base}+stack_{size}-8,
  \qquad RSP_0\bmod16=8.
\]

这是 System V AMD64 函数入口可观察的余数。用户运行时在该槽写零返回地址，
线程入口本身则为 noreturn：函数返回值最终交给 \code{ExitThread}，不能
真的执行 \code{ret} 跳到零。

Kernel 在 irq-save scheduler lock 内取得尚未发布的 Thread 槽，建立动态
Ring 0 双 guard 栈，在其顶部构造统一 UserContext，并初始化干净 FXSAVE64
区。只有运行时元数据、saved frame 和 TLS base 全部就绪后，Thread 才进入
Ready queue。若提前恢复中断，PIT 可以在创建系统调用尚未返回时选择半个
Thread，错误 RIP/RSP 往往表现为“偶发 QEMU 卡住”，本质却是发布协议被破坏。

\begin{failurebox}
父 Thread 不能在 \code{CreateThread} 返回后才把 TID 写入子 TLS。新 Thread
一经 Ready 就可能先运行；子入口必须调用 \code{GetThreadId} 取得自己的 TID。
线程创建给出的是并发关系，不是“父返回必然先于子入口”的顺序保证。
\end{failurebox}

\topic{FS-base：长模式为什么仍保留一小块分段}
8086 用 segment 与 offset 扩展地址表达；保护模式又让 selector 指向带权限
的描述符。x86-64 长模式把普通代码与数据段基址基本固定为零，却保留 FS 和
GS 的独立 64 位 base。历史兼容机制因此成为现代系统放置 TLS 与 per-CPU
数据的实用通道。

本项目明确分工：

\begin{itemize}
  \item GS、\code{SWAPGS} 与 \code{IA32\_KERNEL\_GS\_BASE} 服务 CpuLocal
        和可信特权入口；
  \item FS 与 \code{IA32\_FS\_BASE} 服务 Ring 3 ThreadLocalStorage；
  \item 每次激活 Thread 都与 CR3、TSS.RSP0、FXSAVE 一起恢复 FS-base。
\end{itemize}

TLS 页首保存 \code{ThreadRuntimeState* self}。同一条
\code{mov reg,qword ptr fs:[0]} 指令的线性地址为

\[
  linear = IA32\_FS\_BASE_{current\ thread}+0.
\]

两个 Thread 的机器码和 offset 完全相同，MSR 中不同的 base 使它们取得不同
对象。运行时再检查 self 指针和 64 位 validation tag，避免把任意映射误认
为有效线程状态。

传统 32 位特权入口常执行 \code{mov fs,kernel\_data\_selector}。若机械保留
这条指令，第一次系统调用可能把隐藏 FS base 重置为描述符的零基址。Kernel
自身不使用 FS，因此 SYSCALL、IRET 和调度汇编路径不再重载 FS；
真正切换 Thread 时由架构包装写 \code{IA32\_FS\_BASE} 并读回验证。

\topic{为什么仅靠自旋不能形成通用用户 Mutex}
原子 compare-exchange 已能保证“只有一个 Thread 把 0 改为 1”，却不能决定
竞争失败者应怎样等待。在单 vCPU QEMU 中，若持锁者刚被抢占，竞争者持续
自旋会白白耗尽整个时间片；CPU 显示繁忙，受保护工作却没有推进。

实现把路径分成两层：

\begin{center}
\begin{osgridtabular}{L{0.23\linewidth}L{0.29\linewidth}L{0.36\linewidth}}
路径 & 操作 & 成本与语义 \\ \hline
无竞争 Lock & CAS \(0\rightarrow1\), acquire & 不进入 Kernel \\ \hline
竞争 Lock & wait(word, expected \(=1\)) & Thread 进入 WaitQueue \\ \hline
Unlock & release store 0，wake one & 发布临界区写入并允许一个竞争者重试 \\ \hline
\end{osgridtabular}
\end{center}

acquire/release 不是“让编译器慢一点”的提示。成功 acquire 后的读写不能
移动到取得锁之前；release 之前的临界区写入必须在另一个 Thread 的 acquire
之后可见。futex 只安排睡眠，不替代该内存顺序，也不保存 Mutex owner。

\topic{compare-and-block 封闭丢失唤醒窗口}
错误实现把检查与阻塞分开：

\begin{pdfdiagram}[node distance=9mm and 17mm]
  \node[os state] (read) {waiter 读到 1};
  \node[os hardware,right=of read] (unlock) {owner 写 0\\wake 发现无人};
  \node[os state,right=of unlock] (sleep) {waiter 才入队\\永久睡眠};
  \draw[os arrow] (read) -- node[above]{竞态窗} (unlock);
  \draw[os arrow] (unlock) -- node[above]{通知已消失} (sleep);
\end{pdfdiagram}
\diagramnote{值已经允许前进，但 waiter 错过唯一一次通知；增加 timeout 不能修复协议。}

正确 futex wait 在 scheduler lock 内完成最后一次用户字读取、值比较、取得
entry、登记 WaitQueue 和 Thread 状态切换。wake 也持同一锁。于是全序只可能
有两种有效结果：

\begin{enumerate}[label=\arabic*.]
  \item wake 先提交，wait 读到值已变化并返回 ValueChanged；
  \item wait 先提交 Blocked，wake 随后一定能在 queue 找到它。
\end{enumerate}

Kernel 在锁外先做一次用户地址验证和读取，使合法 demand fault 可以正常处理；
锁内第二次复制才决定是否真正睡眠。若锁内复制失败，操作返回明确状态，不能
在持 scheduler lock 时递归处理页故障。

\topic{private futex 的 key 为什么不是虚拟地址或物理页}
两个 Process 可以同时把 \code{0x60001000} 映射为不同 frame。只用虚拟地址
会让 A 的 unlock 唤醒 B。改用当前物理地址仍不正确：匿名 reservation 在
首次 fault 前没有 frame，fork COW 会暂时共享 frame，写故障又会替换 frame。

因此本阶段使用

\[
  key = (AddressSpaceId,\;AlignDown_4(user\ virtual\ address)).
\]

AddressSpaceId 在地址空间建立时单调取得，不等于 CR3 frame，也不等于可复用
Process 槽。futex word 固定为 \code{uint32\_t}，ABI 对齐为四字节，不依赖
宿主的 \code{unsigned long} 位宽。

\code{PrivateFutexManager} 有 512 个 entry。第一个 waiter 为 key 取得空槽
并初始化 WaitQueue；同 key 后续 waiter 复用。最后一个 waiter 被唤醒或取消
后，\code{ReleaseIfEmpty} 重置 queue 再归还 entry。仍有 waiter 时释放必须
返回 WaitersRemain，容量耗尽则返回 LimitExceeded；实现不能退回无界自旋。

\topic{ConditionVariable 保存变化序列而不是业务条件}
“队列非空”“工作已完成”之类谓词属于调用者数据，并受 Mutex 保护。条件变量
只保存一个 32 位 sequence：

\begin{center}
\begin{osgridtabular}{L{0.43\linewidth}L{0.43\linewidth}}
等待者 & 通知者 \\ \hline
持 Mutex 读取 sequence & 修改受保护谓词 \\ \hline
释放 Mutex & sequence 加一，release \\ \hline
futex wait(sequence, old) & wake one 或 wake all \\ \hline
重新取得 Mutex & 返回继续工作 \\ \hline
循环复查谓词 & 不替等待者判定谓词 \\ \hline
\end{osgridtabular}
\end{center}

通知若落在“释放 Mutex”与“进入内核 wait”之间，内核二次比较看到序列已改变，
等待者不会睡眠。序列允许自然回绕；正确性来自循环复查谓词，而不是把每次
通知当作不可丢失的消息队列。

\topic{Once 用三态区分执行中与已发布}
一个 bool 只能表示“未完成/完成”，无法区分“无人开始”和“某个 Thread 正在
构造共享对象”。本阶段固定：

\[
  0=NotStarted,\qquad 1=Running,\qquad 2=Completed.
\]

唯一 CAS \(0\rightarrow1\) 成功者执行函数；完成后 release store 2 并
wake all。观察到 1 的 Thread 在状态字上等待，观察到 2 的 Thread 以 acquire
直接返回。这样初始化函数在完成前的全部写入对后来调用者可见。

本阶段没有 C++ exception、Thread cancellation 或 robust owner-death。
初始化函数若永不返回，其他 Thread 继续等待。这不是未记录的偶然行为，而是
当前先返回明确错误，等待以后扩展这些失败语义。

\topic{ExitThread、Join 与 ProcessExit 是三种回收边界}
普通 \code{ExitThread} 只终止当前执行流：

\begin{enumerate}[label=\arabic*.]
  \item 保存 64 位 exit value，状态转为 Exited；
  \item 唤醒以目标 TID 为条件的 join waiter；
  \item 切换到另一条 Thread，不关闭 Process 共享对象。
\end{enumerate}

join 为目标记录唯一 join owner。目标退出后，Kernel 在安全点释放其
KernelStack 并 reap 调度槽，把 TID 与 exit value 复制给 joiner；用户运行时
随后撤销目标用户栈和 TLS 页。若先撤用户映射，Kernel 的 saved frame 或退出
路径可能仍引用旧栈；若两个 joiner 都能消费，则同一槽会被重复回收。

\begin{pdfdiagram}[node distance=7mm and 9mm]
  \node[os state] (running) {Running};
  \node[os block,right=of running] (exited) {Exited\\保存 exit value};
  \node[os state,right=of exited] (kernelreap) {Join 成功\\回收 KernelStack/槽};
  \node[os hardware,below=of kernelreap] (userreap) {用户撤映射\\stack + TLS};
  \draw[os arrow] (running) -- node[above]{ExitThread} (exited);
  \draw[os arrow] (exited) -- node[above]{唯一 owner} (kernelreap);
  \draw[os arrow] (kernelreap) -- (userreap);
\end{pdfdiagram}
\diagramnote{逻辑退出、内核安全回收和用户地址撤销是三个顺序明确的时刻。}

最后一条 Thread 退出时，Process 已没有其他执行者，路径提升为 ProcessExit。
显式 ProcessExit、用户异常和某些 fatal 生命周期都会取消 private futex，
从各 WaitQueue 移除 Blocked sibling，终止 Ready/Running sibling，最后关闭
FileTable、FsContext 与 AddressSpace。取消也必须增加 queue、Thread 和
scheduler 的 wake 账本，否则对象虽不再可达，资源快照仍会暴露计数不守恒。

\topic{多线程 exec 为什么先验证候选、再停止 sibling}
exec 替换整个 AddressSpace。兄弟 Thread 的 RIP、RSP 和 FS-base 全部指向
旧映像，不能只让调用者切换 CR3 后继续放任 sibling。另一方面，路径不存在、
ELF 截断、W+X、参数越界或内存不足都属于可恢复错误，不应先杀死旧 Thread。

提交顺序因此为：

\begin{enumerate}[label=\arabic*.]
  \item 在独立候选 AddressSpace 中读取、验证并映射 ELF；
  \item 构造候选参数栈，短暂激活页表验证硬件可用性；
  \item 回到旧地址空间，取消其 private futex wait；
  \item 终止并回收全部 sibling；
  \item 提交调用 Thread 的 CR3、RIP、RSP，并把 FS-base 清零；
  \item 销毁旧地址空间。
\end{enumerate}

步骤一、二失败时旧 Process 完整继续。步骤三以后若内部不变量破坏，旧线程集
已经无法诚实恢复，Kernel 应停机暴露缺陷，不能向用户返回一个伪装的 errno。
fork 的边界不同：只复制调用 Thread，child 不继承其他 Thread 的栈、TLS 或
等待关系。

\topic{munmap 必须认识活动 Thread 和 futex waiter}
若 Thread A 在地址 \(X\) 上等待，Thread B 直接撤销 \(X\)，Kernel queue 会
保存一个不再属于原对象的 key；同地址重新映射后甚至可能接收旧 wake。因此
成功 unmap 会按 AddressSpaceId 与半开区间取消相交 futex，waiter 得到
Cancelled 后回到用户态重新判断。

活动 Thread 的用户栈和 TLS 也不能被其他 Thread 撤销。下次函数调用、IRQ
返回或 \code{fs:[0]} 会立即访问这些区域。Kernel 为每个活动 Thread 记录
stack/TLS 半开区间，unmap 在修改 VMA 前完成重叠预检。只有 join 已在 Kernel
reap 目标后，用户运行时才有权撤销对应映射。

\topic{容量不是一个孤立数字}
三档 QEMU 运行同一实现，只改变资源配置和探针目标：

\begin{center}
\begin{osgridtabular}{L{0.25\linewidth}L{0.20\linewidth}L{0.23\linewidth}L{0.22\linewidth}}
档位 & RAM & Thread 规格 & thread probe \\ \hline
bootstrap & 64 MiB & 每 Process 1 & 明确验证不能再创建 \\ \hline
functional & 256 MiB & 系统 128、每 Process 32 & 总计 32 条同时 Ready \\ \hline
capacity & 64 GiB & 系统 512、每 Process 64 & 总计 64 条并验证上限 \\ \hline
\end{osgridtabular}
\end{center}

每条用户 Thread 还消耗 \(64\,\mathrm{KiB}\) 用户栈 reservation、4 KiB
TLS、一组页表/VMA 描述符、动态 KernelStack 和 FXSAVE 区。逻辑 reservation
不等于全部物理驻留：未触及的匿名栈页不分配 frame；guard 永远没有 VMA。
容量测试既检查“能创建目标数量”，也检查 join 后活动 Thread、栈槽、futex
entry、VMA、页表和 frame 回到基线。

\topic{测试怎样分别证明局部算法与真实硬件路径}
宿主单元测试检查相同 key 复用、不同 AddressSpaceId 隔离、非空 entry 拒绝
释放、容量耗尽、唯一 join owner，以及 Process 终止同时处理 Running、
Ready 和 Blocked sibling。固定种子随机测试在 \(8\times8\) 个 futex key
上执行十万步 acquire、enqueue、wake、cancel、release；独立模型逐步比较
状态和统计。

ABI 集成测试检查系统调用 47--53、48 字节创建请求、16 字节 join 结果和
所有固定宽度常量。真实 QEMU probe 再建立 32 或 64 条 Thread，让它们：

\begin{itemize}
  \item 在同一 Mutex 下共同递增计数；
  \item 用 ConditionVariable 形成等待屏障并执行 notify-all；
  \item 竞争同一个 Once，最终执行次数必须精确为一；
  \item 多次被 PIT 抢占后仍读回各自 FS-base TLS 与 TID；
  \item 返回不同 64 位 exit value，由主 Thread 逐一 join；
  \item 最终输出 TLS、futex 与 join 三类具名证据。
\end{itemize}

测试总数不是接口规格。测试层次、随机可复现、QEMU 总截止/静默截止、
64 GiB 容量门禁，以及成功摘要中的创建、退出、等待、唤醒和回收守恒才是
需要长期保留的约束。

\topic{按控制流走读源码}\par
\begin{enumerate}[label=\arabic*.]
  \item 从 \filepath{source/abi/include/os/abi/thread.hpp} 读固定布局；
  \item 读 \filepath{source/user/src/thread.cpp} 的 guard/TLS/入口/Join；
  \item 读 \filepath{source/user/src/synchronization.cpp} 的用户原子快路径；
  \item 在 \filepath{source/kernel/src/user/system_calls.cpp} 追 ABI copy；
  \item 在 \filepath{source/kernel/src/process/process_runtime.cpp} 找创建提交点、
        compare-and-block、exec 与 exit 组合事务；
  \item 在 \filepath{source/kernel/src/sync/private_futex.cpp} 核对 key 和 entry
        生命周期；
  \item 在 \filepath{source/kernel/src/process/thread_scheduler.cpp} 核对状态
        转换与 sibling 终止；
  \item 最后读 \filepath{source/kernel/src/arch/architecture.asm}，确认系统
        调用与调度返回没有重载 FS。
\end{enumerate}
